%% --- Question --- %%
\chapter*{Interim Project Progress Report Requirement}

%% Overview
\section*{Overview}
For this assignment, you will need to submit both the code that you have produced for your project to date and written responses to the report prompts in the ‘My Assignment’ tab.\\
The aim of this assignment is to make sure that your project meets the requirements for the main submission at the end of the module and that you have made sufficient progress so far.
The assessment of your project will also be based on your use of design and planning techniques and the quality of your writing.\\
The report prompts in the ‘My Assignment’ tab allow you to include images, headings and text formatting.
Make sure you use the features to highlight your points clearly.\\
\textbf{You must not write more than the stated word limit for each prompt or you will lose marks (10\% discretion will be applied, so for a 400 word limit, 440 words or less will not be penalised).
You can write fewer words without any penalty.}\\
It is acceptable to use bullet points to briefly explain your project; for example, the extensions you are adding to the template.
However, you should also include the coding techniques you will use and how complex you believe the extension to be.\\
Don’t worry about your code being a work in progress; we are not looking at the code's stability or quality in this assignment.
However, we want to see evidence that you are making progress and have considered you initial implementation ideas practically.
\section*{Questions}

%% Question 1
\subsection*{Question 1}
Outline your project and which extensions you are building on to the template. (400 words)\\
Your answer should include:
\begin{itemize}
  \item Which of the templates you are extending and why have you chosen to do it.
  \item What extensions have you chosen to do. You should include in your answer:
    \begin{itemize}
      \item any complex coding techniques you will need to use,
(such as arrays of objects, constructor functions, nested looping)
        \item the complexity of the extension
        \item and any expected challenges you will have implementing it.
    \end{itemize}
\end{itemize}

%% Question 2
\subsection*{Question 2}
Discuss the progress you have made on your project so far. (400 words)\\
Your answer should include:
\begin{itemize}
  \item What design work and research you have done while developing your extensions.
    \begin{itemize}
  \item You can include diagrams here in support of your answer. \\This might include but not be limited to: brainstorming, interface designs, code diagrams, and notes on the ideas and code.
  \item You don’t have to include everything you have done up to this point but try and be selective and select the best examples of your process.
\end{itemize}
  \item How much of the code have you written so far.
  \item What do you intend to do next?
\end{itemize}

%% Question 3
\subsection*{Question 3}
Discuss how you will organise you time for the rest of the project. (200 words)
\begin{itemize}
  \item Your answer should be based around a Gantt chart or other time organisation process.
  \begin{itemize}
    \item You can easily make this in a spreadsheet program and export an image or take a screenshot for inclusion here.
    \item Any words on the diagram do not count towards the word count.
    \item If you use a Gantt chart should include the total time of the projects in weeks.
    \item You will also need to divide up the parts of your program and submission into activities for allocating time to on the chart
  \end{itemize}
  \item Your timescales should include progress to this point as well as afterwards.
  \item You should discuss how you have decided on the time you have allocated to each activity and why you have allocated these amounts of time.
  \item You should make reference to your chart to clarify the points you make here.
\end{itemize}

%% Question 4
\subsection*{Question 4}
List any external sources that you have actively utilised in your project. This should include:
\begin{itemize}
  \item any code you have used from external sources directly
  \item any code that you have taken inspiration from but adapted and refined for the project (such as pseudocode algorithms or code pens)
  \item any online help forums you have taken code from (i.e. StackOverflow or library documentation)
  \item any third-party libraries you are using.
\end{itemize}
You do not need to include everything you have read or that has helped you.
Only where you have used or adapted code that appears in your project.

%% Question 5
\subsection*{Question 5}
Upload your work in progress project files as a zip. Please don’t use any other compression file formats.

%% Question 6
\subsection*{Question 6}
Please upload your code in text format (not screenshots) in a PDF file here.
Please clearly label with start and end comments exactly which sections of code you personally wrote without assistance.
5\% of the marks for this coursework are reserved for this part. 

%% Question 7
\subsection*{Question 7}
Discuss how your project will be original and unique (100 words). For example:
\begin{itemize}
  \item What novel features have you developed for your application design.
  \item How you have applied complex coding techniques to your ideas and designs.
  \item How have you adapted and evolved inspiration you have seen from other sources.
  \item Can you justify your ideas in terms of a use case or use group where existing solutions aren't available.
\end{itemize}

%% Review Criteria
\section*{Review Criteria}
In general we will be marking and commenting on the following criteria:
\begin{enumerate}
\item Is your project sufficiently complex enough to meet the criteria of the assignment? - 20\%
  \begin{itemize}
    \item This will be important in determining if you need to add additional extensions or revise them to use more coding techniques.
    \item The marker will indicate in the comments if they think you are taking too much on and should maybe reduce the scope of the project. 
  \end{itemize}
\item Have you made satisfactory progress on the project so far? - 30\% 
  \begin{itemize}
    \item Use this as a guide for how you have progressed so far and what you still need to achieve.
    \item A low score here will indicate that you are in danger of not completing the project in time or with sufficient difficulty.
  \end{itemize}
\item Have you provided a detailed enough plan of how you are going to complete your project? - 20\%
\item Does the design work you have produced (notes, interface designs, code design drawings) effectively detail the development process you have undertaken? - 15\%
\item Has your report been well enough written and accurately communicates your work at the required level of detail? This includes spelling and grammar. - 15\%
\end{enumerate}
